% Chapter 4, typeset from lectures/L04: the README is the opener and the closing
% review, and each appendix is a section, in the same order and with the same subsections.

\chapter{Kernel II: Communication and Application Architecture}
\label{c4:ch:ipc}\label{c4:ch}

{\large\itshape How the parts of an application talk to each other.\par}

\begin{chapterbox}{In this chapter}
\begin{itemize}
  \item Kernel objects for synchronization: semaphores, mutexes and events.
  \item Priority inversion: how it happens, and how a mutex’s priority inheritance bounds it.
  \item Kernel objects for passing data: message queues, FIFOs, and \code{k_poll()} for waiting on several.
  \item Blocking, timeouts and \code{K_NO_WAIT}: what a thread should do when a queue is full.
  \item Memory without \code{malloc()}: static allocation, \code{k_heap} and memory slabs.
  \item zbus: channels, validators, listeners, subscribers and message subscribers.
  \item The State Machine Framework: states, transitions, entry, run and exit actions, and parent states.
  \item Structuring an application so that its parts know about channels and events, not about each other.
\end{itemize}
\end{chapterbox}

A system of threads is only as good as what passes between them. This chapter is the catalogue:
semaphores, mutexes, events, message queues and polling, what each one guarantees, and the trouble
each one causes when it is used for the wrong job. Priority inversion gets a section of its own,
because it is the one bug in this chapter that a correct-looking program can have for years.

The second half is architecture rather than primitives: zbus, which lets the parts of an application
publish to a channel instead of knowing about each other, and the State Machine Framework, which is
the difference between a device's behaviour written down and the same behaviour spread across four
threads. The examples are in \file{lectures/L04/examples}.

% Section 4.1, from lectures/L04/appendix/a_kernel_objects.md.

\section{Kernel Objects: Synchronization, Data and Memory}
\label{c4:sec:kernelobjectssynch}\label{c4:app:a}
Chapter~\ref{c3:ch} was about threads on their own: who runs, when, and for how long. This section
is about threads together: how they wait for each other, protect what they share, hand each other
data, and get memory without \code{malloc()}. By the end you should be able to explain every line
the \file{lectures/L04/examples/msgq} and \file{lectures/L04/examples/priority_inversion} examples
print. \secref{c4:app:b} builds on it with two ways of structuring an application.

\subsection{Waiting, and how long for}
Every kernel call that can block takes a \textbf{timeout} as its last argument, of type
\code{k_timeout_t}:

\begin{booktable}{@{}LL@{}}
  \thead{Timeout} & \thead{Meaning} \\
  \midrule
  \code{K_NO_WAIT} & Do not wait. If the call cannot succeed at once, it fails at once. \\
  \code{K_MSEC(n)}, \code{K_SECONDS(n)}, … & Wait at most that long, then fail. \\
  \code{K_FOREVER} & Wait for as long as it takes. \\
\end{booktable}

A call that fails tells you why through its return value, and the two failures mean different
things. \code{-EBUSY} (or \code{-ENOMSG}, or \code{-ENOMEM}, depending on the object) says \emph{it
was not possible right now, and you asked not to wait}. \code{-EAGAIN} says \emph{you waited as long
as you allowed, and it still did not become possible}. Both are ordinary results, not errors to
ignore: a thread that takes a lock with a timeout and does not check whether it got it is a data
race waiting to happen.

\textbf{An interrupt handler must always pass \code{K_NO_WAIT}}, because an ISR has no thread to put
to sleep. The API documentation marks every function an ISR may call with \emph{isr-ok}, and the
ones that take a timeout add \emph{timeout must be set to \code{K_NO_WAIT} if called from ISR}.

\textbf{A thread waiting on an object waits in a queue ordered by priority.} When the object becomes
available, the highest-priority waiter gets it, and among waiters of equal priority, the one that
has waited longest. A thread that becomes ready this way preempts the current thread if it has a
higher priority, exactly as a thread waking from \code{k_sleep()} would.

\subsection{Semaphores}
A \code{k_sem} is a counter with a limit:

\begin{ccode}
K_SEM_DEFINE(data_ready, 0, 1);   /* Initial count 0, limit 1. */

/* In an interrupt handler, or another thread: */
k_sem_give(&data_ready);

/* In the thread that processes the data: */
if (k_sem_take(&data_ready, K_MSEC(100)) == 0) {
	process();
}
\end{ccode}

\noindent \code{k_sem_take()} decrements the count if it is above zero, and otherwise waits. It
returns \code{0} when it got the semaphore, \code{-EBUSY} when it did not and was told not to wait,
and \code{-EAGAIN} when its timeout expired. \code{k_sem_give()} wakes the first waiter if there is
one, and otherwise increments the count, \textbf{but never above the limit}: three gives to a
semaphore with limit 1 and no waiters leave its count at 1. It never blocks, and may be called from
an ISR.

Semaphores are the tool for \textbf{signalling and counting}: \emph{data is ready}, \emph{an event
happened}, \emph{there are N free buffers}. A semaphore with limit 1 signals that something happened
at least once since the last take, not how many times.

\textbf{A semaphore has no owner.} Any thread may give it, whether or not that thread took it, and
the kernel does not know or care which thread ``holds'' it. That is exactly what makes a semaphore
right for signalling, and wrong as a lock, for a reason the next two sections are about.

\subsection{Mutexes}
A \code{k_mutex} is a lock with an owner:

\begin{ccode}
K_MUTEX_DEFINE(bus_lock);

k_mutex_lock(&bus_lock, K_FOREVER);
/* ... use the bus: the critical section ... */
k_mutex_unlock(&bus_lock);
\end{ccode}

\noindent Three things follow from having an owner:

\begin{itemize}
  \item \textbf{Only the owner may unlock it.} \code{k_mutex_unlock()} from another thread returns
    \code{-EPERM}, and on a mutex that is not locked at all, \code{-EINVAL}.
  \item \textbf{It is recursive.} The owner may lock it again without blocking: the kernel counts the locks,
    and the mutex is released when the owner has unlocked it as many times as it locked it. A helper
    that locks the bus can therefore be called from code that already holds the bus lock.
  \item \textbf{It cannot be used in an ISR.} An interrupt handler is not a thread, so it cannot own
    anything. An ISR that needs to share data with a thread uses a spinlock, \code{irq_lock()}, or
    better, hands the data over through a message queue.
\end{itemize}

And one thing that matters more than all three: \textbf{priority inheritance}.

\subsubsection{Priority inversion}
Three threads share one CPU: \code{high}, \code{medium} and \code{low}. \code{low} locks a resource
that \code{high} also needs. Then:

\begin{enumerate}
  \item \code{low} takes the lock and starts its critical section.
  \item \code{high} becomes ready, preempts \code{low}, tries to take the lock, and blocks.
  \item \code{medium} becomes ready. It needs no lock at all. It has a higher priority than \code{low}, so
    it preempts \code{low}, and runs for as long as it has work.
  \item Only when \code{medium} is done does \code{low} get the CPU back, finish its critical section, and
    release the lock. Only then can \code{high} run.
\end{enumerate}

\code{high}, the most important thread in the system, has waited for \code{medium}, a less important
thread it has nothing to do with. That is \textbf{priority inversion}, and what makes it dangerous
is step 3: \code{high}’s wait is no longer bounded by \code{low}’s critical section, which the
designer knows and keeps short, but by the run time of every thread with a priority between the two,
which nobody bounded. The Mars Pathfinder lander reset itself repeatedly in 1997 for exactly this
reason.

The \file{lectures/L04/examples/priority_inversion} example produces it on purpose. With a semaphore
as the lock, \code{high waited 579 ms}: 69 ms of \code{low}’s remaining work plus all 500 ms of
\code{medium}’s.

\subsubsection{Priority inheritance}
When a thread blocks on a locked \code{k_mutex}, the kernel raises the owner’s priority to the
waiter’s, if the waiter’s is higher. The owner now runs at \code{high}’s priority, so \code{medium}
can no longer preempt it; it finishes its critical section and unlocks, and at the unlock the kernel
restores the priority it had when it locked the mutex. With a mutex in the example, \code{high
waited 79 ms}, just the rest of \code{low}’s critical section, and the trace shows \code{low
(priority 5) gives the mutex back}: \code{low} was running at \code{high}’s priority.

\bookfigure{priority_inversion}{The two runs of the \code{priority_inversion} example on one
scale. A solid bar is a thread with the CPU; a dashed one is a thread that is ready or blocked. With
the semaphore, \code{high} waits for work that has nothing to do with the lock; with the mutex, it
waits only for \code{low}'s critical section.}{c4:fig:inversion}

If the waiting thread gives up because its timeout expires, the kernel lowers the owner’s priority
again, to that of the highest-priority thread still waiting, or back to its own.

What priority inheritance does \textbf{not} do:

\begin{itemize}
  \item \textbf{It does not remove the wait.} \code{high} still waits for \code{low}’s critical section.
    Inheritance bounds the wait by the critical section’s length; keeping that short is still the
    designer’s job.
  \item \textbf{It does not cope well with several mutexes.} The owner’s original priority is saved in each
    mutex when it is locked. If a thread holds several mutexes and releases them in a different order
    than its priority was raised, its priority can be restored wrongly. Zephyr’s documentation
    recommends that a thread lock only one mutex at a time when several are shared between threads of
    different priorities.
  \item \textbf{It does not prevent deadlock.} Two threads that lock two mutexes in opposite orders still
    wait for each other forever. Always take nested locks in the same order.
\end{itemize}

\code{CONFIG_PRIORITY_CEILING} caps how high inheritance may raise a priority. Its default, -128,
allows any.

\subsubsection{On \code{native_sim}: making a thread use CPU time}
Priority inversion needs a thread that \emph{computes}: holds the CPU without blocking. On
\code{native_sim}, code runs in zero simulated time, so a loop of instructions takes no time at all,
and \code{medium} would finish before anything could preempt it. The examples and exercises use
\textbf{\code{k_busy_wait()}} instead, which advances simulated time while keeping the CPU, and
during which interrupts are still handled: a thread that becomes ready meanwhile preempts the busy
thread, as it would on a microcontroller. The examples busy-wait one millisecond at a time, because
a busy wait ends at a fixed deadline even if the thread was preempted in between: one call of
\code{k_busy_wait(100000)} would count time spent preempted as work done.

\subsection{Events}
A \code{k_event} is a set of 32 bits that threads can wait on:

\begin{ccode}
K_EVENT_DEFINE(heater_events);

#define EVENT_START BIT(0)
#define EVENT_FAULT BIT(4)

/* From anywhere, including an ISR: */
k_event_post(&heater_events, EVENT_FAULT);

/* In the thread: wait for any of the events, and clear the ones returned. */
uint32_t events = k_event_wait_safe(&heater_events, EVENT_START | EVENT_FAULT,
				   false, K_FOREVER);
\end{ccode}

\noindent \code{k_event_post()} sets bits, \code{k_event_set()} replaces them all, and
\code{k_event_clear()} clears some. \code{k_event_wait()} returns as soon as any of the requested
bits is set, and \code{k_event_wait_all()} when all of them are. Both return the matching bits, or 0
on a timeout. The \code{_safe} versions also clear the bits they return in the same step, so that an
event posted between the wait returning and a separate \code{k_event_clear()} cannot be lost. Events
need \code{CONFIG_EVENTS=y}.

\textbf{An event object is a set of bits, not a queue.} Posting a bit that is already set does
nothing, so three posts of \code{EVENT_START} before the thread runs are one \code{EVENT_START}, and
the order in which different bits were posted is lost. The \file{lectures/L04/examples/smf} example
shows it happening. Use events for \emph{which conditions are true now}, and a message queue when
the count or the order matters.

\subsection{Message queues}
A \code{k_msgq} passes fixed-size messages between threads, by copying them:

\begin{ccode}
struct sample {
	uint32_t seq;
	int64_t timestamp_ms;
};

/* Room for four samples, allocated statically.
 * The last argument is the alignment. */
K_MSGQ_DEFINE(samples, sizeof(struct sample), 4, alignof(struct sample));

/* The producer: */
struct sample s = {.seq = 1, .timestamp_ms = k_uptime_get()};
int err = k_msgq_put(&samples, &s, K_NO_WAIT);

/* The consumer: */
struct sample s;
k_msgq_get(&samples, &s, K_FOREVER);
\end{ccode}

\noindent \code{k_msgq_put()} copies the message into the queue, and \code{k_msgq_get()} copies it
out, first in, first out. Both return \code{0} on success, \code{-ENOMSG} when they could not and
were told not to wait, and \code{-EAGAIN} when their timeout expired. Put may be called from an ISR
with \code{K_NO_WAIT}, which makes a message queue the standard way to hand data from an interrupt
handler to a thread.

Because messages are copied, \textbf{the sender’s variable can go out of scope as soon as the put
returns}, and the queue’s memory is the only memory involved: \code{max_msgs × msg_size} bytes,
allocated by \code{K_MSGQ_DEFINE()} at link time. Keep messages small, a few tens of bytes. Larger
data travels better as a pointer in a small message, or through a FIFO.

When a thread is already waiting in \code{k_msgq_get()}, a put copies the message straight into that
thread and makes it ready, and if it has the higher priority, it runs before the put even returns.
That is why the \file{lectures/L04/examples/msgq} example prints \code{consumed \#1} before
\code{produced \#1}, and a queue that is \code{0/4} right after the first put.

Other calls: \code{k_msgq_peek()} reads the next message without removing it,
\code{k_msgq_put_front()} puts a message at the head, \code{k_msgq_purge()} discards everything and
wakes blocked senders with \code{-ENOMSG}, and \code{k_msgq_num_used_get()} and
\code{k_msgq_num_free_get()} report how full the queue is.

\subsubsection{When the queue is full}
A full queue is a design decision, not an accident, and the producer’s timeout is where it is made:

\begin{itemize}
  \item \textbf{\code{K_FOREVER}: slow the producer down.} Nothing is lost, and the producer runs at the
    consumer’s pace. Right when every message matters and the producer can wait, such as commands.
  \item \textbf{\code{K_NO_WAIT}: drop the new message.} The producer keeps its rate, and must decide what
    to do with the failure: count it, report it. The only option in an ISR.
  \item \textbf{Make room: drop the oldest.} With one producer, \code{k_msgq_get()} with \code{K_NO_WAIT}
    and a second put keep the newest data, which is usually what sensor data wants.
  \item \textbf{A timeout: wait a little.} Absorbs short bursts without letting a stuck consumer stop the
    producer.
\end{itemize}

The example’s two phases show the first two: with \code{K_NO_WAIT} samples \#8, \#10 and \#12 are
dropped, and with \code{K_FOREVER} nothing is, but from \#21 on every put blocks for 150 ms.

\subsection{FIFOs and LIFOs}
A \code{k_fifo} passes data items \textbf{by reference}:

\begin{ccode}
struct packet {
	void *reserved;      /* The first word belongs to the kernel. */
	size_t length;
	uint8_t data[128];
};

K_FIFO_DEFINE(rx_packets);

/* Put never blocks and may be called from an ISR; get returns NULL on timeout. */
k_fifo_put(&rx_packets, packet);
struct packet *p = k_fifo_get(&rx_packets, K_FOREVER);
\end{ccode}

\noindent Nothing is copied: the FIFO links the items together through their \textbf{first word},
which the kernel overwrites, so every item type starts with a pointer-sized member the code never
touches. The item must stay valid until it is taken out, so it cannot be a local variable of the
sender; it comes from a static pool, a memory slab or a heap. A FIFO has no size limit, because the
items bring their own memory, and it never blocks on put. A \code{k_lifo} is the same, last in,
first out.

Use a message queue for small values and a FIFO for large buffers that already live somewhere.

\subsection{Waiting for several things: \code{k_poll()}}
A thread that must react to a command \emph{or} new data cannot block in \code{k_msgq_get()} on one
of them. \code{k_poll()} waits for any of several conditions:

\begin{ccode}
struct k_poll_event events[] = {
	K_POLL_EVENT_INITIALIZER(K_POLL_TYPE_SEM_AVAILABLE,
				 K_POLL_MODE_NOTIFY_ONLY, &data_ready),
	K_POLL_EVENT_INITIALIZER(K_POLL_TYPE_MSGQ_DATA_AVAILABLE,
				 K_POLL_MODE_NOTIFY_ONLY, &commands),
};

while (true) {
	k_poll(events, ARRAY_SIZE(events), K_FOREVER);

	if (events[0].state == K_POLL_STATE_SEM_AVAILABLE) {
		k_sem_take(&data_ready, K_NO_WAIT);
		process_data();
	}
	if (events[1].state == K_POLL_STATE_MSGQ_DATA_AVAILABLE) {
		k_msgq_get(&commands, &command, K_NO_WAIT);
		handle(&command);
	}
	events[0].state = K_POLL_STATE_NOT_READY;
	events[1].state = K_POLL_STATE_NOT_READY;
}
\end{ccode}

\noindent Three rules, all easy to miss:

\begin{itemize}
  \item \textbf{\code{k_poll()} does not take anything.} It only reports that the semaphore was available:
    the thread still takes it, with \code{K_NO_WAIT}, and must handle failing to, since a thread blocked
    in \code{k_sem_take()} has precedence over a polling one.
  \item \textbf{Check every event’s \code{state}}, since several may be ready at once.
  \item \textbf{Reset every \code{state} to \code{K_POLL_STATE_NOT_READY}} before polling again.
\end{itemize}

Polling needs \code{CONFIG_POLL=y}. Besides semaphores and message queues, it can wait on FIFOs and
on a \code{k_poll_signal}, a lightweight flag raised with \code{k_poll_signal_raise()}.

\subsection{Choosing}
\begin{booktable}{@{}LL@{}}
  \thead{You need to} & \thead{Use} \\
  \midrule
  Protect data shared between threads & \code{k_mutex} \\
  Protect data shared with an ISR & A spinlock or \code{irq_lock()}, held briefly; or pass the data, below \\
  Signal that something happened, or count available resources & \code{k_sem} \\
  Wait for any or all of several conditions & \code{k_event}, or \code{k_poll()} \\
  Pass small values, copied, from a thread or an ISR & \code{k_msgq} \\
  Pass large buffers by reference & \code{k_fifo} \\
  Decouple the parts of an application & zbus, in \secref{c4:app:b} \\
\end{booktable}

\textbf{Never use a semaphore as a lock between threads of different priorities.} It works, until a
thread of medium priority arrives.

\subsection{Memory without malloc}
Every object in this section was defined with a \code{K_..._DEFINE()} macro, at file scope, and so
were the threads and their stacks. That is not a coincidence: \textbf{Zephyr applications allocate
statically}, so that the linker knows all the memory the system will ever use, and the question
\emph{can we run out of memory at run time?} is answered before the firmware ships. For the same
reason, \textbf{Zephyr has no system heap by default}: \code{CONFIG_HEAP_MEM_POOL_SIZE} is 0, and a
call to \code{k_malloc()} does not even link:

\begin{console}
undefined reference to `k_malloc'
\end{console}

\noindent When an application does need memory at run time, Zephyr offers two allocators, each
backed by a static region of its own, so that one part of the system running out cannot starve
another.

\subsubsection{Memory slabs}
A \code{k_mem_slab} hands out fixed-size blocks:

\begin{ccode}
/* Four blocks of 64 bytes, aligned to 8. */
K_MEM_SLAB_DEFINE(packet_slab, 64, 4, 8);

void *block;

if (k_mem_slab_alloc(&packet_slab, &block, K_NO_WAIT) == 0) {
	/* ... */
	k_mem_slab_free(&packet_slab, block);
}
\end{ccode}

\noindent Allocation and freeing take constant time, and since every block has the same size, a slab
cannot fragment: after any sequence of allocations and frees, a free block is always usable. When
the slab is empty, \code{k_mem_slab_alloc()} returns \code{-ENOMEM} without waiting, or waits for a
block to be freed. Allocation may be done in an ISR with \code{K_NO_WAIT}. Slabs are the right tool
for many objects of one type: packets, messages, buffers of a fixed size.

\subsubsection{Heaps}
A \code{k_heap} allocates blocks of any size from its own region:

\begin{ccode}
K_HEAP_DEFINE(frame_heap, 1024);

/* NULL if it could not. */
void *frame = k_heap_alloc(&frame_heap, length, K_MSEC(10));
k_heap_free(&frame_heap, frame);
\end{ccode}

\noindent A heap can fragment, so it can fail even when enough bytes are free in total. Size it with
margin, and use it where sizes genuinely vary. \code{k_malloc()} and \code{k_free()} allocate from
the system heap, which exists only when \code{CONFIG_HEAP_MEM_POOL_SIZE} gives it a size.

The C library’s \code{malloc()} is a separate allocator again. On \code{native_sim} it is the host’s
own \code{malloc()}, which never runs out in practice, so an application that uses it will behave
very differently on a board. That is one more reason to stay with the kernel’s allocators.

\subsection{Summary}
\begin{itemize}
  \item Blocking calls take a timeout. \code{K_NO_WAIT} fails at once, a timeout fails with \code{-EAGAIN},
    and an ISR may only use \code{K_NO_WAIT}.
  \item Waiters are served highest priority first, and a woken thread of higher priority runs at once.
  \item A semaphore is a counter without an owner: use it for signalling and counting.
  \item A mutex has an owner, is recursive, cannot be used in an ISR, and inherits priority. Use it for
    locking, keep critical sections short, and hold one mutex at a time.
  \item Priority inversion lets medium-priority threads delay a high-priority one without bound. Priority
    inheritance bounds the delay by the critical section.
  \item On \code{native_sim}, CPU work is simulated with \code{k_busy_wait()}, which lets higher-priority
    threads preempt.
  \item A \code{k_event} holds conditions as bits: posts merge, and order is lost.
  \item A message queue copies small messages, and a FIFO links large items by reference. What a producer
    does with a full queue is a design decision.
  \item \code{k_poll()} waits for several objects, but takes none of them.
  \item Allocate statically. There is no system heap unless one is configured; memory slabs and heaps give
    run-time allocation from bounded regions.
\end{itemize}

% Section 4.2, from lectures/L04/appendix/b_zbus_and_smf.md.

\section{Application Architecture: zbus and the State Machine Framework}
\label{c4:sec:applicationarchite}\label{c4:app:b}
\secref{c4:app:a} gave threads the means to cooperate. This section is about keeping that
cooperation manageable as an application grows: \textbf{zbus}, which lets the parts of an
application exchange messages without knowing about each other, and the \textbf{State Machine
Framework}, which turns \emph{what should happen now?} into a table rather than a tangle of flags.
By the end you should be able to explain every line the \file{lectures/L04/examples/zbus} and
\file{lectures/L04/examples/smf} examples print.

\subsection{Why decouple}
A temperature sensor thread that calls \code{display_update()}, \code{alarm_check()} and
\code{logger_write()} knows about three other modules, links against them, and has to change
whenever a fourth one wants the temperature. Give it a message queue per consumer instead and it
still knows about all of them, only through their queues. Either way, testing the sensor means
building everything it talks to.

The alternative is to make the \emph{data} the interface. The sensor publishes temperatures to a
\textbf{channel}; whoever needs temperatures observes that channel. The sensor does not know who is
listening, and the consumers do not know who publishes. Adding a consumer changes no existing code,
and a test can publish fake temperatures to exercise the consumers, or observe the channel to check
the sensor. zbus is Zephyr’s implementation of that idea.

\subsection{zbus}
zbus needs \code{CONFIG_ZBUS=y}. It has two kinds of thing: \textbf{channels}, which hold messages,
and \textbf{observers}, which are told when a channel changes.

\subsubsection{Channels}
A channel is a statically allocated message of one type, a lock, and a list of observers:

\begin{ccode}
struct temperature_msg {
	uint32_t seq;
	int32_t millidegrees;
};

ZBUS_CHAN_DEFINE(temperature_chan,              /* The channel's name. */
		 struct temperature_msg,        /* The message type. */
		 temperature_is_valid,          /* A validator, or NULL. */
		 NULL,                          /* User data, or NULL. */
		 ZBUS_OBSERVERS(alarm, display), /* Who is notified, in this order. */
		 ZBUS_MSG_INIT(0));             /* The message's initial value. */
\end{ccode}

\noindent The channel \textbf{owns} its message: there is exactly one copy of it, inside the
channel, and it always holds the last value published. Everything else gets copies, or a look while
the channel is locked.

A \textbf{validator} is a function that checks a message before it is accepted:

\begin{ccode}
static bool temperature_is_valid(const void *msg, size_t msg_size)
{
	const struct temperature_msg *t = msg;

	return (t->millidegrees >= -40000) && (t->millidegrees <= 125000);
}
\end{ccode}

\noindent A message the validator rejects is never stored, no observer hears of it, and the
publisher gets \code{-ENOMSG}. That keeps the check in one place, next to the data, rather than in
every consumer.

\subsubsection{Publishing and reading}
\begin{ccode}
struct temperature_msg t = {.seq = 1, .millidegrees = 21500};
int err = zbus_chan_pub(&temperature_chan, &t, K_MSEC(100));
\end{ccode}

\noindent \code{zbus_chan_pub()} does five things, in order:

\begin{enumerate}
  \item Runs the validator, and returns \code{-ENOMSG} if it rejects the message.
  \item Locks the channel, waiting at most the timeout, and returns \code{-EBUSY} or \code{-EAGAIN} if it
    cannot.
  \item Copies the message into the channel.
  \item Notifies every observer, in the order \code{ZBUS_OBSERVERS()} lists them.
  \item Unlocks the channel.
\end{enumerate}

If notifying an observer fails, for example because a subscriber’s queue is full,
\code{zbus_chan_pub()} returns that error, \textbf{although the message was stored}, and the
remaining observers are still notified. The one exception is running out of memory for a message
subscriber’s copy, which stops the notifications there.

Reading copies the message out, under the lock:

\begin{ccode}
struct temperature_msg t;
zbus_chan_read(&temperature_chan, &t, K_MSEC(100));
\end{ccode}

\noindent \code{zbus_chan_claim()} and \code{zbus_chan_finish()} lock the channel and give access to
the message in place, for large messages that should not be copied; hold the claim briefly, since it
blocks every publisher and reader.

\textbf{A note on error numbers.} The \file{lectures/L04/examples/zbus} example compares the return
value with \code{-ENOMSG} rather than printing it, because the number depends on the C library.
\code{native_sim} uses the host’s, where \code{ENOMSG} is 42; Zephyr’s own minimal C library defines
it as 35. Always compare with the name.

\subsubsection{Observers}
There are four kinds of observer. What distinguishes them is \textbf{which thread runs their code},
and \textbf{what they receive}:

\begin{booktable}{@{}LLL@{}}
  \thead{Observer} & \thead{Runs in} & \thead{Receives} \\
  \midrule
  Listener \newline \code{ZBUS_LISTENER_DEFINE(name, callback)} & The publisher's thread, inside
    \code{zbus_chan_pub()}, with the channel locked & A call; it reads the message in place with
    \code{zbus_chan_const_msg()} \\
  \addlinespace
  Subscriber \newline \code{ZBUS_SUBSCRIBER_DEFINE(name, queue_size)} & A thread of its own,
    waiting in \code{zbus_sub_wait()} & A notification that the channel changed; it reads the
    channel itself \\
  \addlinespace
  Message subscriber \newline \code{ZBUS_MSG_SUBSCRIBER_DEFINE(name)} & A thread of its own,
    waiting in \code{zbus_sub_wait_msg()} & A copy of every message \\
  \addlinespace
  Asynchronous listener \newline \code{ZBUS_ASYNC_LISTENER_DEFINE(name, callback)} & A work queue
    & A copy of the message, passed to the callback \\
\end{booktable}

\textbf{A listener is synchronous.} Its callback runs before \code{zbus_chan_pub()} returns, in the
publisher’s thread, with the channel locked, so it is effectively part of the publisher. In the
example, \code{alarm saw \#1 ... in thread 'sensor_thread'} is printed before \code{zbus_chan_pub()
returned 0}. A listener must therefore be short and must never block, just like an interrupt
handler, and it cannot publish to its own channel, whose lock its thread already holds. A listener
is the right tool for something quick and urgent: setting an alarm output, waking a thread.

\textbf{A subscriber is told, not given.} Its notification is a pointer to the channel, put into a
message queue of \code{queue_size} entries. The subscriber’s thread reads the channel when it gets
to run, and sees whatever the channel holds \emph{then}:

\begin{ccode}
ZBUS_SUBSCRIBER_DEFINE(display, 4);

static void display_main(void)
{
	const struct zbus_channel *chan;
	struct temperature_msg t;

	while (zbus_sub_wait(&display, &chan, K_FOREVER) == 0) {
		zbus_chan_read(chan, &t, K_FOREVER);
		show(&t);
	}
}
\end{ccode}

\noindent \textbf{So a subscriber can miss messages.} If two messages are published before the
subscriber’s thread runs, both notifications are waiting, and both reads return the second message.
In the example, the display sees \#3 twice and never sees \#2. That is not a bug: for a display,
only the latest value matters, and skipping stale ones is exactly right. When a subscriber’s queue
is full, the publisher waits for room for as long as its timeout allows, \textbf{still holding the
channel’s lock}, so the subscriber cannot read the channel to catch up; when the timeout expires,
the notification is lost and \code{zbus_chan_pub()} returns the error. Size the queue for the bursts
the subscriber must survive.

\textbf{A message subscriber gets everything.} zbus copies each message into a buffer (a
\code{net_buf}, from the heap by default) and queues the copy, and \code{zbus_sub_wait_msg()}
returns it:

\begin{ccode}
ZBUS_MSG_SUBSCRIBER_DEFINE(logger);

while (zbus_sub_wait_msg(&logger, &chan, &t, K_FOREVER) == 0) {
	log_reading(&t);
}
\end{ccode}

\noindent In the example, the logger sees both \#2 and \#3. The cost is a buffer per queued message,
and it needs \code{CONFIG_ZBUS_MSG_SUBSCRIBER=y}. Use it when every message matters, as for a logger
or a protocol.

\textbf{Asynchronous listeners} need \code{CONFIG_ZBUS_ASYNC_LISTENER=y}: their callback runs later,
from a work queue, with a copy of the message. They suit work too slow for a listener that does not
justify a thread of its own.

\subsubsection{Priorities decide the order}
A listener’s output always appears inside the publish. Everything else happens when the scheduler
says so. In the example, the sensor thread has the highest priority, so the subscribers only run
once it sleeps; give the sensor a lower priority than the subscribers, and each subscriber runs as
soon as its notification is queued, before \code{zbus_chan_pub()} has even returned, and then the
display sees \#2 too.

\code{CONFIG_ZBUS_PRIORITY_BOOST}, enabled by default, can raise a publisher’s priority while it
holds a channel, above the highest-priority observer thread, so that observers cannot preempt a
publish halfway. It only knows an observer’s priority once the observer’s thread has called
\code{zbus_obs_attach_to_thread()}; the examples do not, so no boost happens there.

\subsubsection{Designing with channels}
\begin{itemize}
  \item \textbf{One channel per kind of data}, named for the data (\code{temperature_chan}), not for its
    producer or consumer.
  \item \textbf{Small, flat messages.} Every publish and read copies them.
  \item \textbf{Choose the observer by what it needs:} a listener for urgent and quick, a subscriber for
    \emph{the latest value}, a message subscriber for \emph{every value}.
  \item \textbf{Validate in the channel}, not in the consumers.
\end{itemize}

\subsection{The State Machine Framework}
Most devices are state machines, whether or not their code says so. A heater is idle, heating,
holding a temperature, or in a fault, and what an event means depends on which: \emph{temperature
reached} matters while heating, and means nothing while idle. Written as flags and \code{if}
statements, that knowledge is spread over every function that handles an event. Written as a state
machine, it is in one place, and each state’s code only handles what can happen in that state.

The State Machine Framework (SMF) needs \code{CONFIG_SMF=y}, and
\code{CONFIG_SMF_ANCESTOR_SUPPORT=y} for parent states.

\subsubsection{States}
A state is a \code{struct smf_state} with up to three \textbf{actions}, functions that take the
state machine’s object:

\begin{itemize}
  \item \textbf{entry}, run when the state is entered;
  \item \textbf{run}, run each time the state machine runs while in the state, typically once per event;
  \item \textbf{exit}, run when the state is left.
\end{itemize}

The whole state machine is an array of states:

\begin{ccode}
enum heater_state { OPERATIONAL, IDLE, HEATING, HOLD, FAULT };

static const struct smf_state states[] = {
    [OPERATIONAL] = SMF_CREATE_STATE(operational_entry, operational_run,
                                     operational_exit, NULL, NULL),
    [IDLE] = SMF_CREATE_STATE(idle_entry, idle_run, NULL,
                              &states[OPERATIONAL], NULL),
    [HEATING] = SMF_CREATE_STATE(heating_entry, heating_run, heating_exit,
                                 &states[OPERATIONAL], NULL),
    [HOLD] = SMF_CREATE_STATE(hold_entry, hold_run, NULL,
                              &states[OPERATIONAL], NULL),
    [FAULT] = SMF_CREATE_STATE(fault_entry, fault_run, fault_exit, NULL, NULL),
};
\end{ccode}

\noindent \code{SMF_CREATE_STATE(entry, run, exit, parent, initial)} takes the three actions, any of
which may be \code{NULL}, the state’s parent, and an initial child state, used only with
\code{CONFIG_SMF_INITIAL_TRANSITION}.

\subsubsection{The object and its context}
The state machine works on an object of the application’s own, whose \textbf{first member} must be a
\code{struct smf_ctx}, so that a pointer to the object is also a pointer to the context:

\begin{ccode}
struct heater {
	struct smf_ctx ctx;   /* Must be first. */
	uint32_t event;       /* What the run actions look at. */
};
\end{ccode}

\noindent \code{SMF_CTX(&heater)} casts the object to its context for SMF’s functions, and every action receives
the object as a \code{void *}.

\subsubsection{Running it}
Three functions do everything:

\begin{itemize}
  \item \textbf{\code{smf_set_initial(ctx, state)}} starts the state machine in a state, running the entry
    actions of the state and its ancestors, outermost first.
  \item \textbf{\code{smf_run_state(ctx)}} runs the current state’s run action once.
  \item \textbf{\code{smf_set_state(ctx, state)}}, called from an action, transitions to another state: it
    runs the exit actions of the states being left and the entry actions of the states being entered.
\end{itemize}

A typical application runs the state machine once per event, from a thread:

\begin{ccode}
while (true) {
	k_msgq_get(&heater_events, &heater.event, K_FOREVER);
	smf_run_state(SMF_CTX(&heater));
}
\end{ccode}

\noindent A run action decides what the event means in its state:

\begin{ccode}
static enum smf_state_result heating_run(void *obj)
{
	struct heater *heater = obj;

	switch (heater->event) {
	case EVENT_TEMP_REACHED:
		smf_set_state(SMF_CTX(heater), &states[HOLD]);
		return SMF_EVENT_HANDLED;
	case EVENT_STOP:
		smf_set_state(SMF_CTX(heater), &states[IDLE]);
		return SMF_EVENT_HANDLED;
	default:
		return SMF_EVENT_PROPAGATE;
	}
}
\end{ccode}

\noindent \code{smf_set_state()} must not be called from an exit action, where the state machine is
already in the middle of a transition; SMF logs an error and ignores it.
\code{smf_set_terminate(ctx, value)} stops the state machine, and makes \code{smf_run_state()}
return the value from then on.

\subsubsection{Parent states}
With \code{CONFIG_SMF_ANCESTOR_SUPPORT=y}, states form a hierarchy, and two things change.

\textbf{Events propagate upwards.} A run action returns \code{SMF_EVENT_HANDLED} or
\code{SMF_EVENT_PROPAGATE}. If the state handled the event, or transitioned, that is the end of it.
Otherwise the parent’s run action gets the event, then the grandparent’s, until one handles it. A
state without a run action propagates everything. Behaviour that several states share is written
once, in their parent: in the example, a fault in IDLE, HEATING or HOLD is handled by OPERATIONAL
alone.

\textbf{Transitions only leave and enter what changes.} A transition runs exit actions from the
current state up to, but not including, the lowest ancestor that the old and new state have in
common, and entry actions from below that ancestor down to the new state. Moving from HEATING to
HOLD, both children of OPERATIONAL, runs HEATING’s exit and HOLD’s entry, and OPERATIONAL is neither
left nor entered. Moving to FAULT, which is outside OPERATIONAL, runs OPERATIONAL’s exit as well,
and the example’s trace shows \code{OPERATIONAL exit}. Whatever a parent sets up in its entry
action, such as switching a motor on, therefore stays set up while the state machine moves between
its children.

Without \code{CONFIG_SMF_ANCESTOR_SUPPORT}, the parent argument is ignored, every state is on its
own, and a run action’s return value does not matter: that is a \textbf{flat} state machine.

\subsubsection{Driving it with events}
The state machine’s thread needs events to run on, and how they arrive matters:

\begin{itemize}
  \item \textbf{A message queue} delivers every event, in order, and can carry data with it. It is the
    default choice.
  \item \textbf{A \code{k_event}} delivers \emph{which} events happened, as bits. It is simple and
    ISR-friendly, but events of the same kind merge and their order is lost: the
    \file{lectures/L04/examples/smf} example’s README shows \code{START}, \code{STOP}, \code{START},
    \code{START} being handled as \code{START}, \code{START}, \code{STOP}. Use it only when that cannot
    matter.
  \item \textbf{zbus}: the state machine’s thread subscribes to the channels whose messages are its events.
\end{itemize}

\textbf{Keep the state machine itself free of threads.} If its code only reacts to the event it is
given and records what it does, it can be tested by calling it directly with a sequence of events,
with no threads and no timing at all. Exercise Set~\ref{c4:set:4} is built that way.

\subsection{Summary}
\begin{itemize}
  \item zbus decouples an application’s parts: publishers and observers only know channels.
  \item A channel owns one statically allocated message of one type, and a validator can reject messages
    before they are stored.
  \item \code{zbus_chan_pub()} validates, locks, copies, notifies observers in order, and unlocks.
  \item A listener runs inside the publish, in the publisher’s thread, with the channel locked: short, and
    never blocking.
  \item A subscriber is notified and reads the latest value, so it can miss intermediate messages. A message
    subscriber gets a copy of every message.
  \item Thread priorities decide when subscribers run, relative to the publisher.
  \item Compare error codes with their names: \code{native_sim} uses the host’s C library.
  \item SMF describes a state machine as a table of states with entry, run and exit actions.
  \item \code{smf_set_initial()} starts it, \code{smf_run_state()} runs it once per event, and
    \code{smf_set_state()} transitions.
  \item With parent states, unhandled events propagate upwards, and transitions do not leave or re-enter
    common ancestors.
  \item Feed a state machine from a message queue when order and count matter, and keep its logic free of
    threads, so that it can be tested directly.
\end{itemize}

% Section 4.3, from lectures/L04/README.md: the objectives, the evaluation
% questions, and the next lecture.

\section{Review}
\label{c4:sec:review}

\needspace{6\baselineskip}
\subsection*{What you should be able to do now}
After this chapter, you should:

\begin{itemize}
  \item Be able to choose the right kernel object for a synchronization or communication problem.
  \item Understand why a semaphore used as a lock allows unbounded priority inversion, and a mutex does not.
  \item Be able to pass data between threads with a message queue, and handle a full queue deliberately.
  \item Understand how Zephyr applications avoid dynamic memory, and what to use when they cannot.
  \item Be able to decouple an application’s parts with zbus channels, and choose between its kinds of
    observer.
  \item Be able to express a device’s behaviour as a hierarchical state machine with SMF.
\end{itemize}

\needspace{6\baselineskip}
\subsection*{Questions to test yourself}
You should be able to explain:

\begin{enumerate}
  \item What is the difference between a semaphore and a mutex, beyond the count?
  \item What is priority inversion, and why is the medium-priority thread the one that makes it unbounded?
  \item How does priority inheritance bound it, and what does it not solve?
  \item When should a producer use \code{K_NO_WAIT} on a full queue, and when a timeout?
  \item Why do Zephyr applications usually allocate kernel objects and buffers statically?
  \item What is the difference between a zbus listener, a subscriber and a message subscriber, and in which
    thread does each one run?
  \item Why can a zbus subscriber miss a message, and which observer cannot?
  \item What does a parent state in SMF save you from writing?
  \item On \code{native_sim}, why must a thread that simulates CPU work use \code{k_busy_wait()} rather than
    a loop?
\end{enumerate}

\needspace{6\baselineskip}
\subsection*{Looking ahead}
Chapter~\ref{c5:ch} is about writing a driver:

\begin{itemize}
  \item The device driver model: API structs, config and data, instances and init levels.
  \item A binding and a Kconfig option for a new part.
  \item Emulators: hardware that exists only in software, on \code{native_sim}’s I2C bus.
  \item Testing the driver with ztest, and running the tests with Twister.
\end{itemize}

You will write a driver for an I2C temperature sensor that does not exist.

% Section 4.4, from lectures/L04/appendix/c_exercises.md.

\section{Exercises}
\label{c4:sec:exercises}\label{c4:app:c}
These exercises reinforce the concepts from \secref{c4:app:a} and \secref{c4:app:b}.

\begin{aside}
\textbf{How to check your work.} Write each exercise set as its own application under
\file{lectures/L04/exercises}, with the file names given there, and run \code{make test} from the
repository root. A set’s tests switch on as soon as its files exist, and a set you have not started
is reported as \code{SKIP}. Each application grows through its set, and the tests check where it
ends up. Programs run on \code{native_sim} in simulated time, so their output, timestamps included,
is compared with the expected output line by line. The threads’ priorities, start delays and sleeps
given below are part of the exercise: change one, and the output changes.

The solutions are published in the companion repository, under
\file{lectures/L04/appendix/solutions}, a few days after the lecture. Try each exercise yourself
before you read them.
\end{aside}

Remember \code{source env.sh} in every new terminal. Build and run an exercise from the repository
root:

\begin{shell}
west build -p -b native_sim/native/64 lectures/L04/exercises/exercise1
./build/zephyr/zephyr.exe -stop_at=3 -no-rt
\end{shell}

% ---------------------------------------------------------------------
\exerciseset{Readings in a queue}
\label{c4:set:1}
Create an application named \code{exercise1} to work in for this exercise set, with a
\code{CMakeLists.txt}, a \code{prj.conf} and a \code{src/main.c}.

\exercise{A sampler and a slower storage thread}{Code}
\label{c4:ex:1.1}
In this exercise you will connect a fast producer and a slow consumer with a message queue.

\task{Tasks}
\begin{itemize}
  \item In \code{prj.conf}, set one tick per millisecond, as this chapter’s examples do:
    \code{CONFIG_SYS_CLOCK_TICKS_PER_SEC=1000}.
  \item Define a message type:
\begin{ccode}
struct reading {
  uint32_t seq;
  int64_t timestamp_ms;
};
\end{ccode}
  \item Define a message queue named \code{readings} with room for \textbf{3} readings.
  \item Write a thread \code{sampler}, priority \textbf{5}, which takes 8 readings, numbered 1 to 8. For
    each:
\begin{itemize}
  \item Fill in a \code{struct reading} with the number and \code{k_uptime_get()}.
  \item Put it into the queue \textbf{without waiting}. If that fails, print \code{[<t> ms] queue full,
    dropped \#<n>}, with the reading’s number. Otherwise print \code{[<t> ms] sampled \#<n>, queue
    <used>/3}.
  \item Sleep 100 ms.
\end{itemize}
After the eighth reading, print \code{[<t> ms] sampler done}.
  \item Write a thread \code{storage}, priority \textbf{6}, which forever takes a reading from the queue,
    waiting as long as it takes, prints \code{[<t> ms]   stored \#<n>, <age> ms old} (three spaces after
    the timestamp), where the age is the uptime minus the reading’s timestamp, and sleeps 350 ms.
  \item Start both threads with \code{K_THREAD_DEFINE()}, a stack of 1024 bytes, and no start delay.
\end{itemize}

Print every timestamp with \code{\%4lld}. With \code{-stop_at=3}, the output must be:

\begin{console}
[   0 ms] sampled #1, queue 1/3
[   0 ms]   stored #1, 0 ms old
[ 101 ms] sampled #2, queue 1/3
[ 202 ms] sampled #3, queue 2/3
[ 303 ms] sampled #4, queue 3/3
[ 351 ms]   stored #2, 250 ms old
[ 404 ms] sampled #5, queue 3/3
[ 505 ms] queue full, dropped #6
[ 606 ms] queue full, dropped #7
[ 702 ms]   stored #3, 500 ms old
[ 707 ms] sampled #8, queue 3/3
[ 808 ms] sampler done
[1053 ms]   stored #4, 750 ms old
[1404 ms]   stored #5, 1000 ms old
[1755 ms]   stored #8, 1048 ms old
\end{console}

\noindent \textbf{Hint:} \code{K_MSGQ_DEFINE()}’s last argument is the message’s alignment:
\code{alignof(struct reading)}, from \code{<stdalign.h>}. \code{k_msgq_num_used_get()} returns how
many messages are queued.

\task{Reflection}
\begin{itemize}
  \item \code{sampled \#1} is printed before \code{stored \#1}. Why? Give \code{storage} priority 4 instead,
    and run it again: what changes, and why does the queue read \code{0/3}?
  \item Why does the sampler put without waiting, rather than with \code{K_FOREVER}? Which kind of code
    could not wait at all?
\end{itemize}

\exercise{Keep the newest}{Code}
\label{c4:ex:1.2}

\task{Tasks}
Change the sampler so that a full queue drops the \textbf{oldest} reading instead of the new one:

\begin{itemize}
  \item When the put fails, take the oldest reading out of the queue without waiting, print \code{[<t> ms]
    queue full, dropped \#<n>} with \textbf{its} number, and put the new reading again.
  \item Then print \code{sampled \#<n>} for the new reading, as for any other.
\end{itemize}

The output must now be:

\begin{console}
[   0 ms] sampled #1, queue 1/3
[   0 ms]   stored #1, 0 ms old
[ 101 ms] sampled #2, queue 1/3
[ 202 ms] sampled #3, queue 2/3
[ 303 ms] sampled #4, queue 3/3
[ 351 ms]   stored #2, 250 ms old
[ 404 ms] sampled #5, queue 3/3
[ 505 ms] queue full, dropped #3
[ 505 ms] sampled #6, queue 3/3
[ 606 ms] queue full, dropped #4
[ 606 ms] sampled #7, queue 3/3
[ 702 ms]   stored #5, 298 ms old
[ 707 ms] sampled #8, queue 3/3
[ 808 ms] sampler done
[1053 ms]   stored #6, 548 ms old
[1404 ms]   stored #7, 798 ms old
[1755 ms]   stored #8, 1048 ms old
\end{console}

\task{Reflection}
\begin{itemize}
  \item Compare the readings stored, and their ages, with Exercise~\ref{c4:ex:1.1}. Which policy does a
    temperature logger want? Which would a queue of commands want?
  \item Make the queue 8 readings long. Is anything dropped now? How old is the last reading when it is
    stored? Did the bigger queue help?
  \item How much RAM does the queue’s buffer use on \code{native_sim/native/64}, and when is it allocated?
  \item Dropping the oldest takes two calls, a get and a put. Why is that safe in this program, and when
    would it not be?
\end{itemize}

% ---------------------------------------------------------------------
\exerciseset{Priority inversion}
\label{c4:set:2}
Create an application named \code{exercise2} to work in for this exercise set.

\exercise{A log behind a semaphore}{Code}
\label{c4:ex:2.1}
In this exercise you will produce priority inversion, measure it, and fix it.

\task{Tasks}
\begin{itemize}
  \item In \code{prj.conf}, set \code{CONFIG_SYS_CLOCK_TICKS_PER_SEC=1000}.
  \item Create a \code{Kconfig} file with a menu named \code{Application} holding a \code{bool} option
    \code{APP_USE_MUTEX}, with the prompt \code{Protect the log with a k_mutex instead of a k_sem}, and
    end it so that Zephyr’s own options are still available. The option stays unused until
    Exercise~\ref{c4:ex:2.2}.
  \item In \code{src/main.c}, protect a shared log with a binary semaphore named \code{log_lock}, created
    available (initial count 1, limit 1).
  \item Write a function \code{trace(who, what)} that prints \code{[<t> ms] <who> (priority <p>) <what>},
    with the uptime as \code{\%4lld}, \code{who} left-aligned in 7 characters (\code{\%-7s}), and the
    calling thread’s current priority from \code{k_thread_priority_get()}.
  \item Write a function \code{compute(ms)} that uses CPU time by calling \code{k_busy_wait(USEC_PER_MSEC)}
    \code{ms} times.
  \item Write three threads, all with a stack of 1024 bytes:
\begin{itemize}
  \item \code{logger}, priority \textbf{9}, no start delay: traces \code{takes the lock}, takes it, computes
    for \textbf{60} ms, traces \code{gives the lock back}, and gives it back.
  \item \code{control}, priority \textbf{5}, start delay \textbf{10} ms: records the uptime, traces
    \code{needs the lock}, takes it, traces \code{has the lock}, gives it back, and prints \code{control
    waited <ms> ms}, the time since the uptime it recorded.
  \item \code{network}, priority \textbf{7}, start delay \textbf{20} ms: traces \code{starts sending},
    computes for \textbf{300} ms, and traces \code{is done}.
\end{itemize}
\end{itemize}

The output must be:

\begin{console}
[   0 ms] logger  (priority 9) takes the lock
[  11 ms] control (priority 5) needs the lock
[  21 ms] network (priority 7) starts sending
[ 321 ms] network (priority 7) is done
[ 360 ms] logger  (priority 9) gives the lock back
[ 360 ms] control (priority 5) has the lock
control waited 349 ms
\end{console}

\noindent \textbf{Hint:} why \code{compute()} busy-waits one millisecond at a time rather than
calling \code{k_busy_wait(60000)} once is in \secref{c4:app:a}, \emph{On \code{native_sim}: making a
thread use CPU time}.

\task{Reflection}
\begin{itemize}
  \item Before you run it, predict how long \code{control} waits. Then account for the 349 ms, millisecond
    by millisecond: who ran when?
  \item \code{network} never touches the lock. How did it delay \code{control} by 300 ms?
\end{itemize}

\exercise{The mutex}{Code}
\label{c4:ex:2.2}

\task{Tasks}
\begin{itemize}
  \item When \code{CONFIG_APP_USE_MUTEX} is enabled, protect the log with a \code{k_mutex} named
    \code{log_lock} instead of the semaphore, and change nothing else. Hide the difference behind two
    macros, such as \code{lock()} and \code{unlock()}, so that the threads’ code is the same in both
    builds.
  \item Create a file \code{mutex.conf} that enables the option.
\end{itemize}

Built with \code{-DEXTRA_CONF_FILE=mutex.conf}, the output must be:

\begin{console}
[   0 ms] logger  (priority 9) takes the lock
[  11 ms] control (priority 5) needs the lock
[  60 ms] logger  (priority 5) gives the lock back
[  60 ms] control (priority 5) has the lock
control waited 49 ms
[  60 ms] network (priority 7) starts sending
[ 360 ms] network (priority 7) is done
\end{console}

\noindent Built without it, the output must be unchanged from Exercise~\ref{c4:ex:2.1}.

\task{Reflection}
\begin{itemize}
  \item Why is \code{logger}’s priority 5 when it gives the lock back? What is it after the unlock?
  \item \code{network} was ready at 21 ms. Why does it only start sending at 60 ms?
  \item \code{control} still waited 49 ms. Could any lock bring that down to zero? What would?
  \item With the mutex, make \code{control} wait at most 20 ms, with \code{k_mutex_lock(&log_lock, K_MSEC(20))},
and trace \code{gave up} when that fails. What happens to \code{logger}’s priority when
\code{control} gives up, and what does \code{network} do then?
\end{itemize}

% ---------------------------------------------------------------------
\exerciseset{A battery on zbus}
\label{c4:set:3}
Create an application named \code{exercise3} to work in for this exercise set.

\exercise{The channel and a listener}{Code}
\label{c4:ex:3.1}
In this exercise you will publish battery readings on a zbus channel, and react to them.

\task{Tasks}
\begin{itemize}
  \item Enable zbus in \code{prj.conf}.
  \item Define the message type:
\begin{ccode}
struct battery_msg {
  uint8_t percent;
  bool charging;
};
\end{ccode}
  \item Define a channel named \code{battery_chan}, initially \code{\{.percent = 100, .charging = false\}},
    with a validator that rejects any message with \code{percent} above 100.
  \item Define a listener named \code{led}, observing the channel, whose callback prints \code{  led: red}
    (two spaces first) when the battery is below 20 \% and not charging, and \code{  led: green}
    otherwise.
  \item Write a thread \code{gauge}, priority \textbf{6}, stack 1024, no start delay, that publishes these
    readings, in this order:
\begin{ccode}
static const struct battery_msg readings[] = {
  {.percent = 80, .charging = false}, {.percent = 15, .charging = false},
  {.percent = 150, .charging = false}, {.percent = 15, .charging = true},
  {.percent = 60, .charging = true},
};
\end{ccode}

\noindent For each, it prints \code{gauge: <percent> \%}, followed by \code{, charging} if the
battery is charging, publishes the reading with a timeout of 50 ms, prints \code{gauge: rejected} if
the channel rejected it, and sleeps 100 ms.
\end{itemize}

The output must be:

\begin{console}
gauge: 80 %
  led: green
gauge: 15 %
  led: red
gauge: 150 %
gauge: rejected
gauge: 15 %, charging
  led: green
gauge: 60 %, charging
  led: green
\end{console}

\noindent \textbf{Hint:} A listener callback reads the message in place with
\code{zbus_chan_const_msg(chan)}. In a format string, \code{\%\%} prints a percent sign.

\task{Reflection}
\begin{itemize}
  \item In which thread does the \code{led} callback run? Check it: enable \code{CONFIG_THREAD_NAME=y}, and
    print \code{k_thread_name_get(k_current_get())} in the callback.
  \item What number does \code{zbus_chan_pub()} return for the 150 \% reading on \code{native_sim}? Why
    should the code compare it with \code{-ENOMSG} rather than with that number?
\end{itemize}

\exercise{A display that subscribes}{Code}
\label{c4:ex:3.2}

\task{Tasks}
\begin{itemize}
  \item Define a subscriber named \code{display}, with a notification queue of \textbf{4}, and add it to the
    channel’s observers after \code{led}.
  \item Write a thread \code{display_thread}, priority \textbf{7}, stack 1024, no start delay, which forever
    waits for a notification, reads the channel, and prints \code{  display: <percent> \%}, followed by
    \code{, charging} if the battery is charging.
\end{itemize}

The output must now be:

\begin{console}
gauge: 80 %
  led: green
  display: 80 %
gauge: 15 %
  led: red
  display: 15 %
gauge: 150 %
gauge: rejected
gauge: 15 %, charging
  led: green
  display: 15 %, charging
gauge: 60 %, charging
  led: green
  display: 60 %, charging
\end{console}

\task{Reflection}
\begin{itemize}
  \item Remove the gauge’s 100 ms sleep. What does the display print now, and why? Did it receive four
    notifications?
  \item Keep the sleep removed, and give \code{display_thread} priority 5, above the gauge. What does the
    display print now?
  \item Why does the \code{led} never miss a reading, whatever the priorities?
\end{itemize}

% ---------------------------------------------------------------------
\exerciseset{A garage door as a state machine}
\label{c4:set:4}
Create an application named \code{exercise4} to work in for this exercise set, with
\code{target_sources(app PRIVATE src/main.c src/door.c)} in its \code{CMakeLists.txt}.

\exercise{The state machine}{Code}
\label{c4:ex:4.1}
In this exercise you will write a hierarchical state machine with SMF, behind an API that tests can
drive directly.

\task{Tasks}
Copy this header into \code{src/door.h} unchanged. It is the interface the tests use:

\begin{ccode}
#ifndef DOOR_H_
#define DOOR_H_

#include <stdbool.h>

/** What can happen to the door. */
enum door_event {
	DOOR_EVENT_BUTTON,   /**< The button was pressed. */
	DOOR_EVENT_OPENED,   /**< The door reached its top end switch. */
	DOOR_EVENT_CLOSED,   /**< The door reached its bottom end switch. */
	DOOR_EVENT_OBSTACLE, /**< The light barrier saw something in the way. */
};

/** The states the door can be in. */
enum door_state {
	DOOR_STATE_CLOSED,
	DOOR_STATE_OPENING,
	DOOR_STATE_OPEN,
	DOOR_STATE_CLOSING,
	DOOR_STATE_STOPPED,
};

/** @brief Start the state machine, in the closed state, with the motor off
 *         and never started. */
void door_init(void);

/** @brief Run the state machine once, for one event. */
void door_handle(enum door_event event);

/** @brief Return the state the door is in. */
enum door_state door_get_state(void);

/** @brief Return whether the motor is running. */
bool door_motor_on(void);

/** @brief Return how many times the motor has been started since door_init(). */
unsigned int door_motor_starts(void);

#endif /* DOOR_H_ */
\end{ccode}

\noindent In \code{prj.conf}, enable SMF with parent states. In \code{src/door.c}, implement the
header with SMF, with these six states:

\begin{booktable}{@{}llp{9.5em}L@{}}
  \thead{State} & \thead{Parent} & \thead{Entry} & \thead{Handles} \\
  \midrule
  \code{MOVING} & none & Motor on, and count the start & \code{BUTTON}: go to \code{STOPPED}. Every other event: ignore it. Exit: motor off. \\
  \code{CLOSED} & none &  & \code{BUTTON}: go to \code{OPENING}. \\
  \code{OPENING} & \code{MOVING} &  & \code{OPENED}: go to \code{OPEN}. Anything else: pass it to the parent. \\
  \code{OPEN} & none &  & \code{BUTTON}: go to \code{CLOSING}. \\
  \code{CLOSING} & \code{MOVING} &  & \code{CLOSED}: go to \code{CLOSED}. \code{OBSTACLE}: go to \code{OPENING}. Anything else: pass it to the parent. \\
  \code{STOPPED} & none &  & \code{BUTTON}: go to \code{CLOSING}. \\
\end{booktable}

Events a state does not list are ignored. \code{door_get_state()} returns the state the door is in,
and \code{MOVING} is never one: the door is always in one of the five states of the header.

\textbf{Hint:} \code{MOVING}’s entry and exit actions are the only code that touches the motor. Keep
the state in a variable that each leaf state’s entry action sets.

\exercise{Drive it}{Code}
\label{c4:ex:4.2}

\task{Tasks}
In \code{src/main.c}, call \code{door_init()}, print the door’s status as \code{start}, and then
handle these events one at a time, printing the status after each, labelled with the event’s name:

\begin{ccode}
static const enum door_event script[] = {
	DOOR_EVENT_BUTTON,   DOOR_EVENT_OBSTACLE, DOOR_EVENT_OPENED,
	DOOR_EVENT_BUTTON,   DOOR_EVENT_OBSTACLE, DOOR_EVENT_OPENED,
	DOOR_EVENT_BUTTON,   DOOR_EVENT_BUTTON,   DOOR_EVENT_BUTTON,
	DOOR_EVENT_CLOSED,   DOOR_EVENT_OBSTACLE,
};
\end{ccode}

\noindent Print each status with \code{"\%-8s -> \%-7s motor \%-3s (started \%u times)\\n"}: the
label, the state’s name, \code{on} or \code{off}, and the number of motor starts. The output must
be:

\begin{console}
start    -> CLOSED  motor off (started 0 times)
BUTTON   -> OPENING motor on  (started 1 times)
OBSTACLE -> OPENING motor on  (started 1 times)
OPENED   -> OPEN    motor off (started 1 times)
BUTTON   -> CLOSING motor on  (started 2 times)
OBSTACLE -> OPENING motor on  (started 2 times)
OPENED   -> OPEN    motor off (started 2 times)
BUTTON   -> CLOSING motor on  (started 3 times)
BUTTON   -> STOPPED motor off (started 3 times)
BUTTON   -> CLOSING motor on  (started 4 times)
CLOSED   -> CLOSED  motor off (started 4 times)
OBSTACLE -> CLOSED  motor off (started 4 times)
\end{console}

\noindent The tests also drive \code{door.c} directly, without \code{main.c}, through sequences of
events the script does not contain.

\task{Reflection}
\begin{itemize}
  \item \code{OBSTACLE} while closing sends the door back to \code{OPENING}, and the motor start count stays
    at 2. Which rule of SMF makes that happen?
  \item Which state handles \code{OBSTACLE} while the door is opening, and what does it do with it?
  \item Change \code{OPENING}’s run action to return \code{SMF_EVENT_HANDLED} for every event. What happens
    when the button is pressed while the door is opening? Which tests fail?
  \item Remove \code{CONFIG_SMF_ANCESTOR_SUPPORT=y} from \code{prj.conf}. Does it still build? What does the
    program print now, and why?
  \item Why can \code{door.c} be tested without threads, timers, or \code{native_sim}’s clock?
\end{itemize}

